\documentclass[11pt,a4paper]{article}
\usepackage[margin=2.5cm]{geometry}
\usepackage{amsmath,amssymb,amsthm}
\usepackage{xcolor}
\usepackage{booktabs}
\usepackage{enumitem}
\usepackage{hyperref}
\hypersetup{colorlinks=true, linkcolor=blue!60!black, urlcolor=blue!60!black, citecolor=blue!60!black}

\usepackage{titlesec}
\titleformat{\section}{\large\bfseries}{\thesection.}{0.5em}{}
\titleformat{\subsection}{\normalsize\bfseries}{\thesubsection.}{0.5em}{}

\setlength{\parindent}{0pt}
\setlength{\parskip}{6pt}

\title{\textbf{LLMs as Resonant Structures}\\[4pt]
\large From Green's Functions to Nonlinear Spectral Theory}
\author{Robin Langer}
\date{}

\begin{document}
\maketitle

\begin{abstract}
We develop a spectral theory of large language models in three layers. First,
the classical picture: an LLM is a resonant structure whose Green's function
encodes its natural modes---a prompt excites the structure, and the response is
determined by the interaction between forcing function and eigenmodes. Second,
the nonlinear correction: the Inverse Scattering Transform shows that certain
nonlinear systems are ``secretly linear'' via a nonlinear change of basis, and
sparse autoencoders (SAEs) perform exactly this role for LLM activations. Third,
the global obstruction: the local linearisations on competence patches form a
presheaf on the latent manifold, and polysemanticity is the failure of the
gluing axiom---measurable by sheaf cohomology. Each level subsumes the last.
The boundary between ``linear'' and ``nonlinear'' is not a wall. It is a choice
of basis.
\end{abstract}

\section{The Punchline}

``Linear'' versus ``nonlinear'' is not a property of the physics. It is a
property of your coordinate system.

The Korteweg--de Vries equation, governing shallow water waves, looks nonlinear
in physical coordinates. In scattering-data coordinates, it is as linear as the
heat equation. The nonlinearity was never in the waves. It was in the description.

The same is true of large language models. An LLM's internal dynamics are
massively nonlinear---attention, layer normalisation, GeLU activations. But
sparse autoencoders reveal that the activations decompose into interpretable,
approximately independent features. In the feature basis, the representation is
approximately diagonal. The nonlinearity is absorbed into the change of coordinates.

This paper develops the analogy in three stages, each subsuming the last:
\begin{enumerate}[leftmargin=*]
  \item \textbf{Green's functions.} A linear spectral decomposition of the LLM
    as a resonant structure with natural modes. Valid locally, and a good first
    approximation.
  \item \textbf{The Inverse Scattering Transform.} A nonlinear spectral theory
    that explains how SAEs find hidden linearity in a nonlinear system---the same
    trick that makes solitons ``secretly linear.''
  \item \textbf{Sheaf structure.} The local linearisations form a presheaf on
    the latent manifold. Polysemanticity is the failure of the gluing axiom.
    The first cohomology $H^1$ measures the degree of context-dependence.
\end{enumerate}

\section{Green's Functions: The Linear Theory}

\subsection{Impulse responses}

Given a linear differential operator $L$ on a domain $\Omega$, the Green's
function $G(x,y)$ solves
\[
  L\, G(\cdot, y) = \delta(y),
\]
where $\delta(y)$ is the Dirac delta centred at $y$. This is the system's
response to a point impulse: strike the drum at exactly one point and observe
how the entire membrane responds.

The power of the Green's function lies in \emph{superposition}. Any input $f$
decomposes as $f(x) = \int_\Omega f(y)\,\delta(y)\,dy$, so the solution to
$Lu = f$ is
\[
  u(x) = \int_\Omega G(x,y)\, f(y)\, dy.
\]
The operator $L$ is the fixed structure---the drum, the membrane, the resonant
cavity. The Green's function encodes everything about what that structure
\emph{can do}. The input $f$ merely selects which modes get excited.

\subsection{Spectral decomposition}

If $L$ has eigenfunctions $\varphi_n$ with eigenvalues $\lambda_n$, then:
\[
  G(x,y) = \sum_{n} \frac{\varphi_n(x)\,\varphi_n(y)}{\lambda_n}.
\]
Each eigenmode $\varphi_n$ is a natural vibration pattern of the system. A small
eigenvalue $\lambda_n$ means the corresponding mode resonates strongly---it takes
little energy to excite it. The input $f = \sum c_n \varphi_n$ simply sets the
amplitudes: the response is $u = \sum (c_n / \lambda_n)\,\varphi_n$.

\subsection{The LLM as resonant structure}

Calling a large language model a ``next-token predictor'' is technically correct
and completely beside the point. The training objective (next-token prediction)
is the \emph{bow}---the mechanism that shaped the instrument during training.
What emerged is the instrument itself: a resonant structure with natural modes,
preferred patterns of response, directions in activation space that ring when
struck.

A prompt does not ``instruct'' the LLM. It \emph{excites} it---a forcing
function applied to a pre-existing vibrational structure. The response is
determined by the interaction between the forcing function (prompt) and the
natural modes (learned representations).

\medskip
\begin{center}
\renewcommand{\arraystretch}{1.3}
\begin{tabular}{@{}ll@{}}
\toprule
\textbf{PDE / Green's function} & \textbf{LLM} \\
\midrule
Linear operator $L$ & Trained network (weights fixed at inference) \\
Domain $\Omega$ & Token/embedding space \\
Green's function $G(x,y)$ & Response map: prompt $\to$ output distribution \\
Eigenmodes $\varphi_n$ & SAE features / representation directions \\
Eigenvalues $\lambda_n$ & Feature activation thresholds \\
Input $f(x)$ & Prompt \\
Solution $u = G * f$ & Completion \\
Resonance & Prompt strongly activating a coherent feature \\
\bottomrule
\end{tabular}
\end{center}

\medskip
This picture already illuminates several phenomena. Prompt engineering becomes
spectral engineering: constructing a forcing function $f$ that excites the right
modes. Fine-tuning changes $L$ (reshapes the cavity); prompting changes $f$
(excites it differently). RLHF reshapes the cavity to amplify some modes and
damp others---the base model and the chat model are different instruments.

This is the linear story. It is correct as far as it goes. But LLMs are not linear.

\section{The Nonlinearity Objection, Taken Seriously}

Green's functions decompose the response into harmonic modes. But purely
harmonic decompositions have well-known limitations: they struggle with shock
waves, jets, and sharp impulses. They approximate these phenomena, but there
are more efficient descriptions. In the language of signal processing, Fourier
analysis assumes global smoothness that nonlinear systems violate.

The standard response---``the network is locally linear around any operating
point, so we can compute a local Jacobian''---is true but insufficient. Local
linearisation is the content of the Hartman--Grobman theorem, and it works near
equilibria. But the interesting behaviour of LLMs happens far from equilibrium:
when the model is composing concepts, resolving ambiguity, or navigating between
incompatible interpretations.

PCA applied to LLM activations is the linear spectral theory---the Green's
function decomposition. It captures the dominant modes but misses the
nonlinear interactions between them. Sparse autoencoders outperform PCA
precisely because they apply a \emph{nonlinear} decomposition that can
separate features packed into fewer dimensions than there are features.

What we need is not a defence of linearity. We need a nonlinear spectral theory.

\section{Solitons: The Nonlinear Spectral Theory}

\subsection{The Inverse Scattering Transform}

The Korteweg--de Vries (KdV) equation
\[
  u_t + 6uu_x + u_{xxx} = 0
\]
governs shallow water waves. It is nonlinear (the $uu_x$ term). Its soliton
solutions---localised waves that propagate without changing shape and survive
collisions intact---were among the great surprises of 20th-century mathematical
physics.

The surprise deepened when Gardner, Greene, Kruskal, and Miura discovered
the Inverse Scattering Transform (IST) in 1967. The IST solves KdV using
exactly the same pipeline as Green's functions solve linear PDEs:

\medskip
\begin{center}
\renewcommand{\arraystretch}{1.3}
\begin{tabular}{@{}ll@{}}
\toprule
\textbf{Linear PDE (heat equation)} & \textbf{Nonlinear PDE (KdV)} \\
\midrule
Fourier transform the initial data & Compute scattering data of $u(x,0)$ \\
Each mode evolves independently & Each piece evolves independently \\
(multiply by $e^{-k^2 t}$) & (multiply by simple exponentials) \\
Inverse Fourier transform & Inverse scattering \\
\bottomrule
\end{tabular}
\end{center}

\medskip
The scattering data consists of two parts: a \emph{discrete spectrum} (finitely
many eigenvalues $\lambda_1, \ldots, \lambda_N$, each with a norming constant
$c_n$) and a \emph{continuous spectrum} (a reflection coefficient $R(k)$). Each
discrete eigenvalue corresponds to exactly one soliton. Its eigenvalue determines
its speed and amplitude. A 3-soliton solution has exactly 3 discrete modes.

The key fact: the eigenvalues do not change under KdV flow. They are conserved
quantities. Solitons survive collisions not because of some mysterious nonlinear
balancing act, but because eigenvalues are invariants of the evolution. In
scattering-data coordinates, KdV is as linear as the heat equation.

The nonlinearity was never in the physics. It was in the choice of coordinates.

\subsection{SAEs as the Inverse Scattering Transform for LLMs}

The analogy between the IST and sparse autoencoders is structural, not
metaphorical:

\medskip
\begin{center}
\renewcommand{\arraystretch}{1.3}
\begin{tabular}{@{}ll@{}}
\toprule
\textbf{IST (soliton theory)} & \textbf{SAE (LLM interpretability)} \\
\midrule
Wave profile $u(x)$ & Activation vector in residual stream \\
Scattering transform & SAE encoder (nonlinear: ReLU) \\
Discrete eigenvalues (solitons) & Monosemantic features \\
Continuous spectrum (radiation) & Polysemantic residual \\
Inverse scattering & SAE decoder \\
Eigenvalues are invariants & Features persist across layers \\
Nonlinear PDE (KdV) & Nonlinear transformer computation \\
\bottomrule
\end{tabular}
\end{center}

\medskip
In both cases, the essential move is the same: a nonlinear change of basis that
reveals hidden linearity. The scattering transform takes a nonlinear wave and
decomposes it into modes that evolve independently. The SAE encoder takes a
superposition of entangled features and decomposes it into sparse, approximately
independent activations.

This explains why SAEs succeed where PCA fails. PCA is the Fourier transform---a
\emph{linear} change of basis. It works when the modes are orthogonal and the
dynamics are linear. But the superposition hypothesis tells us that LLMs pack
more features than dimensions, so features are non-orthogonal. You need a
\emph{nonlinear} decomposition to untangle them, just as you need the IST
(not the Fourier transform) to decompose a nonlinear wave.

PCA is the Green's function. SAEs are the Inverse Scattering Transform.

\subsection{The three regimes}

Not all nonlinear systems admit the IST. The ones that do are called
\emph{integrable}. The analogy extends to a classification:

\medskip
\begin{center}
\renewcommand{\arraystretch}{1.3}
\begin{tabular}{@{}lll@{}}
\toprule
\textbf{Regime} & \textbf{PDE analogue} & \textbf{LLM analogue} \\
\midrule
Integrable & Solitons (clean discrete spectrum) & Fully monosemantic model \\
Nearly integrable & Perturbed solitons (slow radiation) & Mostly monosemantic, some polysemantic \\
Non-integrable & Turbulence (no clean decomposition) & Heavily polysemantic \\
\bottomrule
\end{tabular}
\end{center}

\medskip
Real LLMs sit in the middle row. SAEs work well enough to extract many clean
features, but the decomposition is never perfect. The residual polysemanticity
is the continuous spectrum---the radiation that the IST separates from the
solitons but cannot eliminate.

\section{The Sheaf Structure}

The IST provides a global linearisation when one exists. But what happens when
the system is only \emph{locally} linearisable---when different regions of input
space admit different clean decompositions that may be mutually incompatible?

This is the situation for real LLMs, and the right language for it is sheaf theory.

\subsection{The presheaf of local linearisations}

Let $M$ denote the LLM's latent manifold (the space of activation vectors,
or equivalently, the space of inputs under the embedding map). Give $M$ its
natural topology: nearby points are inputs that produce similar activations.

Define a presheaf $\mathcal{F}$ on $M$ as follows:
\begin{itemize}[leftmargin=*]
  \item To each open set $U \subseteq M$, assign $\mathcal{F}(U)$ = the set of
    local linearisations on $U$. Concretely, this is a local SAE decomposition:
    a set of feature directions that cleanly decompose all activations in $U$
    into sparse, interpretable components.
  \item For $V \subseteq U$, the restriction map
    $\mathcal{F}(U) \to \mathcal{F}(V)$ restricts the decomposition to the
    smaller region. If a linearisation works on $U$, it works on any $V \subseteq U$.
\end{itemize}

This is always at least a presheaf. The question is whether it satisfies the
sheaf axioms.

\subsection{Polysemanticity is the failure of gluing}

A sheaf requires two conditions beyond being a presheaf:

\textbf{Locality (separation).} If two local linearisations agree on every
piece of an open cover, they are the same linearisation. This holds: if two SAE
decompositions produce identical features on every sub-region, they are the same
decomposition.

\textbf{Gluing.} If local linearisations on overlapping patches $\{U_i\}$
agree on all intersections $U_i \cap U_j$, they glue to a global linearisation
on $\bigcup U_i$. This is where it fails.

A feature that activates as ``financial institution'' on patch $U_1$ (finance
context) and ``river edge'' on patch $U_2$ (geography context) provides two
local sections that cannot be glued into a coherent global section. The word
``bank'' has a clean, monosemantic meaning locally on each patch. But the local
meanings are incompatible---they do not extend to a global feature.

This is polysemanticity, stated sheaf-theoretically: the presheaf
$\mathcal{F}$ fails the gluing axiom. Local linearisations exist but are
globally incompatible.

\subsection{Cohomology measures polysemanticity}

The obstruction to gluing lives in the first \v{C}ech cohomology
$H^1(M, \mathcal{F})$. This is a standard construction: given an open cover
$\{U_i\}$ and compatible local sections on each $U_i$, $H^1$ measures
whether they extend to a global section.

\begin{itemize}[leftmargin=*]
  \item $H^1 = 0$: every compatible family of local linearisations glues to a
    global one. The presheaf is a sheaf. The model is globally monosemantic.
  \item $H^1 \neq 0$: there exist compatible local sections that do not glue.
    The model has irreducible polysemanticity---context-dependence that cannot
    be resolved by choosing better features.
\end{itemize}

The dimension of $H^1$ measures \emph{how polysemantic} the model is: how
many independent ways the local linearisations can fail to be globally coherent.

\subsection{Holonomy and the Laplace--Beltrami connection}

If we equip $M$ with a Riemannian metric (induced, for instance, by the Fisher
information of the output distribution), we can study its differential geometry.
The Laplace--Beltrami operator $\Delta_g$ on $M$ generalises the Laplacian; its
eigenfunctions are the natural vibration modes of the manifold.

On a flat manifold, these reduce to Fourier modes---the Green's function
picture. On a curved manifold, they capture the geometry of the representation
space. Competence at a task corresponds to the manifold being locally flat in
the relevant region: the eigenfunctions locally look like Fourier modes, and
the dynamics are locally linear.

The holonomy of the manifold---what happens when you parallel-transport a
feature vector around a closed loop in input space---measures the global
curvature. If the holonomy is trivial (the vector returns unchanged), the local
frames are globally compatible: the presheaf is a sheaf. If the holonomy is
non-trivial, the vector returns rotated: the same feature means something
different after traversing a loop through different contexts. This is
polysemanticity as parallel transport failure.

A training objective that minimises the holonomy of the Laplace--Beltrami
connection is, in sheaf-theoretic terms, attempting to \emph{sheafify} the
representation: to make local linearisations globally compatible by flattening
the manifold via gradient descent.

\subsection{Should we kill $H^1$?}

There is a tension here worth preserving. The superposition hypothesis says that
LLMs deliberately pack more features than dimensions---features are
non-orthogonal precisely because this is efficient compression. A perfectly
sheafified model (one where every local linearisation glues globally) would need
representations as high-dimensional as the feature space. The polysemanticity
is not a bug; it is a compression strategy.

The right question may not be ``make $H^1 = 0$'' but rather ``understand $H^1$
well enough to navigate it.'' Just as an engineer does not need to eliminate
all nonlinearity from a system---only to understand which nonlinearities matter
and when---an interpretability researcher does not need to eliminate
polysemanticity, only to characterise when and how features change meaning
across contexts.

\section{Case Study: The RSK Transformer}

We now present a concrete case where the spectral theory is directly observable:
a small transformer trained to invert combinatorial bijections~\cite{rsk-transformer}.
In the language developed above, this is an \emph{integrable} system---one where
the presheaf is a sheaf, the decomposition is global, and the full spectral
structure can be recovered.

\subsection{Setup}

The Robinson--Schensted--Knuth (RSK) correspondence is a fundamental bijection
in algebraic combinatorics mapping permutations $\sigma \in S_n$ to pairs
$(P, Q)$ of standard Young tableaux of the same shape. The \emph{RSK transformer}
is a 1.2M-parameter encoder that learns the inverse: given $(P, Q)$,
recover~$\sigma$.

The architecture is a 6-layer transformer encoder ($d = 128$, 8 heads) with
structured 2D token embeddings:
\[
  \mathrm{emb}(v, r, c, t) = E_{\mathrm{val}}(v) + E_{\mathrm{row}}(r)
    + E_{\mathrm{col}}(c) + E_{\mathrm{tab}}(t),
\]
where $v$ is the entry value, $(r, c)$ its position in the tableau, and
$t \in \{P, Q\}$ identifies which tableau it belongs to. The same architecture,
without modification, also learns the Hillman--Grassl and Burge correspondences.

\subsection{The drum shape: 2D embeddings as domain geometry}

The ablation study reveals a spectral geometry result. The 2D structured
embeddings define the geometry of the domain $\Omega$---the shape of the drum:

\medskip
\begin{center}
\renewcommand{\arraystretch}{1.3}
\begin{tabular}{@{}lrl@{}}
\toprule
\textbf{Embedding structure} & \textbf{Accuracy} & \textbf{Spectral interpretation} \\
\midrule
Full 2D (row + col + tableau) & 100\% & Correct drum shape $\Rightarrow$ correct eigenmodes \\
1D sequential position & 72\% & Wrong geometry $\Rightarrow$ wrong modes, some survive \\
No tableau identity ($P$ vs $Q$) & 6\% & Domain halved $\Rightarrow$ lost half the spectrum \\
No spatial information & 0\% & No domain $\Rightarrow$ no modes $\Rightarrow$ no response \\
\bottomrule
\end{tabular}
\end{center}

\medskip
This is a direct analogue of Kac's question ``Can one hear the shape of a
drum?''~\cite{kac1966}. Removing the 2D structure changes the drum from a
two-dimensional membrane to a one-dimensional string: some modes survive (72\%),
but those encoding the 2D tableau geometry are destroyed. Removing the
$P$-vs-$Q$ identity glues two drums together, making the network unable to
distinguish which half is vibrating.

\subsection{SAE features as eigenmodes}

Sparse autoencoder analysis reveals interpretable features corresponding
precisely to the natural modes $\varphi_n$:

\begin{itemize}[leftmargin=*]
  \item \textbf{Insertion-order detectors}---temporal eigenmodes that activate
    when the network recognises which step of the RSK algorithm it is inverting.
  \item \textbf{Step-specific locators}---spatial eigenmodes that fire for
    specific tableau cell positions $(r, c)$.
  \item \textbf{Burge rule features}---algebraic eigenmodes encoding the local
    recurrence $G(i,j) = \max(G(i{-}1,j),\, G(i,j{-}1)) - G(i{-}1,j{-}1) + a_{ij}$.
\end{itemize}

These are not abstract directions requiring post hoc interpretation---they are
recognisable steps of known algorithms, recovered as the natural basis of the
network's internal representation.

\subsection{Algorithmic generalisation as resonance}

Training on only 0.00004\% of $S_{15}$ yields 99.99\% exact-match accuracy on
held-out data. In the Green's function picture, this is expected: you do not
need to strike a drum at every point to determine its eigenmodes. A sparse set
of impulses suffices to identify the spectral structure, because each
eigenfunction is defined over the entire domain. Once the modes are
learned, any input decomposes into the same basis.

The MLP baseline collapses from 95.67\% at $n = 10$ to 3.07\% at $n = 15$:
an MLP has no spatial structure to support extended eigenmodes. It is a lookup
table pretending to be a drum.

\subsection{The RSK transformer as an integrable system}

In the framework developed above, the RSK transformer sits in the
\emph{integrable} regime. The evidence:
\begin{itemize}[leftmargin=*]
  \item The SAE decomposition is global: the same features work across the
    entire input space ($S_{15}$). There is no polysemanticity---each feature
    has one meaning everywhere. The presheaf is a sheaf. $H^1 = 0$.
  \item 100\% accuracy means the spectral decomposition is complete: no
    information is lost in the radiation/residual.
  \item The ablation study shows what happens when you \emph{break} the
    sheaf structure: removing geometric information (the domain topology)
    introduces curvature into the representation, forces polysemanticity,
    and accuracy degrades.
  \item The MLP baseline is a non-integrable system: no spatial structure
    to support coherent modes, high effective $H^1$, poor generalisation.
\end{itemize}

The RSK transformer offers a laboratory for the spectral theory of neural
networks precisely because it is small enough to be fully interpretable yet
rich enough to learn deep algebraic structure---and integrable enough that the
full spectral picture is accessible.

\section{Conclusion: The Three Levels}

The spectral theory of LLMs has three levels, each answering a question the
previous cannot:

\medskip
\begin{center}
\renewcommand{\arraystretch}{1.3}
\begin{tabular}{@{}llll@{}}
\toprule
\textbf{Level} & \textbf{Mathematics} & \textbf{LLM tool} & \textbf{What it explains} \\
\midrule
1. Linear & Green's functions & PCA & Natural modes exist \\
2. Nonlinear & Inverse Scattering & SAEs & Hidden linearity in nonlinear systems \\
3. Global & Sheaf cohomology & $H^1$ analysis & Why local decompositions fail globally \\
\bottomrule
\end{tabular}
\end{center}

\medskip
At Level~1, the LLM is a drum. Prompts excite its natural frequencies. This is
true, useful, and incomplete.

At Level~2, the drum is nonlinear---but ``nonlinear'' is a property of the
coordinates, not the physics. SAEs find the scattering-data basis in which the
representation becomes diagonal. Solitons survive collisions because eigenvalues
are invariants; monosemantic features persist across layers for the same reason.

At Level~3, we confront the global question. Local linearisations exist on
competence patches, but they may not glue. Polysemanticity is the cohomological
obstruction. The dimension of $H^1$ measures how context-dependent the model's
representations truly are---and whether that context-dependence is a defect to
be repaired or a compression strategy to be understood.

An LLM is not a machine that follows instructions. It is a medium with a
physics---and that physics has a spectral theory.

\begin{thebibliography}{9}
\bibitem{rsk-transformer}
R.~Langer.
\emph{RSK Transformer: Learning inverse combinatorial bijections}.
\url{https://huggingface.co/RobBobin/rsk-transformer}, 2025.

\bibitem{kac1966}
M.~Kac.
Can one hear the shape of a drum?
\emph{The American Mathematical Monthly}, 73(4):1--23, 1966.

\end{thebibliography}

\end{document}
